\documentclass[letterpaper, 10 pt, conference]{ieeeconf}
\IEEEoverridecommandlockouts
\overrideIEEEmargins

\usepackage{amsmath,amssymb,amsfonts}
\usepackage{graphicx,listings}
\usepackage{cite}

\title{\bf The Trust Kernel: A Distributed Architecture for Governed Autonomy}
\author{[Authors anonymized for submission]}

\begin{document}
\maketitle
\thispagestyle{empty}
\pagestyle{empty}

%%%%%%%%%%%%%%%%%%%%%%%%%%%%%%%%%%%%%%%%%%%%%%%%%%%%%%%%%%%%%%%%%%%%%%%%%%%%%%%%
\begin{abstract}
We present the Trust Kernel, a distributed architecture that operationalizes the GCAT and Rigel frameworks across networked nodes. Each node maintains a local GCAT state and enforces admissibility prior to execution, while distributed attestation ensures global consistency under adversarial conditions. A recovery layer based on projected Rigel dynamics enables coordinated stabilization across nodes. We establish safety under Byzantine faults using standard BFT assumptions and define a distributed recoverability metric. The architecture demonstrates how capacity-bounded autonomy and recoverability can be enforced in a decentralized system.
\end{abstract}

%%%%%%%%%%%%%%%%%%%%%%%%%%%%%%%%%%%%%%%%%%%%%%%%%%%%%%%%%%%%%%%%%%%%%%%%%%%%%%%%
\section{Introduction}

Prior work defined admissibility and recoverability as geometric properties of system state. However, real systems are distributed: decisions are made across multiple nodes with partial information and potential adversarial behavior.

This paper introduces the Trust Kernel, a distributed architecture that enforces admissibility locally while maintaining global consistency through attestation and consensus. The key challenge is to ensure that node-level enforcement composes into system-wide guarantees.

%%%%%%%%%%%%%%%%%%%%%%%%%%%%%%%%%%%%%%%%%%%%%%%%%%%%%%%%%%%%%%%%%%%%%%%%%%%%%%%%
\section{System Model}

We consider a network of $n$ nodes, each maintaining a local state:
\begin{equation}
\mathbf{x}_i = (g_i, c_i, a_i, t_i)
\end{equation}

Each node computes its admissibility margin:
\begin{equation}
m(\mathbf{x}_i) = \Lambda(\mathbf{x}_i) - a_i
\end{equation}

Nodes communicate over a partially synchronous network.

%%%%%%%%%%%%%%%%%%%%%%%%%%%%%%%%%%%%%%%%%%%%%%%%%%%%%%%%%%%%%%%%%%%%%%%%%%%%%%%%
\section{Architecture}

The Trust Kernel consists of three layers:

\subsection{Validation Layer}

Each node enforces:
\begin{equation}
m(\mathbf{x}_i) \ge 0
\end{equation}

prior to executing actions. This ensures local admissibility.

\subsection{Attestation Layer}

Nodes exchange signed state messages:

\begin{verbatim}
STEGVERSE_MSG {
    node_id: UUID,
    timestamp: u64,
    gcat_state: [f64; 4],
    signature: Ed25519,
    proof: MerklePath
}
\end{verbatim}

Messages are aggregated into a consensus state $\mathbf{X}$.

Trust edges are defined as:
\begin{equation}
t_{ij} = \min(t_i, t_j)
\end{equation}

\subsection{Recovery Layer}

Nodes apply distributed Rigel dynamics:

\begin{equation}
u_i = -K_i(\tilde{\mathbf{x}}_i - \mathbf{R}) + \sum_{j \in N(i)} w_{ij} (\tilde{\mathbf{x}}_j - \tilde{\mathbf{x}}_i)
\end{equation}

Projection onto the simplex ensures:
\begin{equation}
\sum_k \tilde{x}_{i,k} = 1
\end{equation}

%%%%%%%%%%%%%%%%%%%%%%%%%%%%%%%%%%%%%%%%%%%%%%%%%%%%%%%%%%%%%%%%%%%%%%%%%%%%%%%%
\section{Consensus and Byzantine Model}

\begin{assumption}
At most $f < n/3$ nodes are Byzantine.
\end{assumption}

Nodes run a BFT consensus protocol (e.g., PBFT or HotStuff) to agree on $\mathbf{X}$.

\begin{theorem}[Safety]
If $f < n/3$, all honest nodes agree on the same committed state.
\end{theorem}

\textit{Argument.}
This follows from standard BFT safety guarantees under partial synchrony.

%%%%%%%%%%%%%%%%%%%%%%%%%%%%%%%%%%%%%%%%%%%%%%%%%%%%%%%%%%%%%%%%%%%%%%%%%%%%%%%%
\section{Compositional Safety}

\begin{proposition}
If all honest nodes enforce local admissibility and consensus ensures consistent state views, then system-wide admissibility violations can be detected and blocked.
\end{proposition}

\textit{Interpretation.}
Local enforcement combined with global agreement prevents inconsistent execution.

%%%%%%%%%%%%%%%%%%%%%%%%%%%%%%%%%%%%%%%%%%%%%%%%%%%%%%%%%%%%%%%%%%%%%%%%%%%%%%%%
\section{Distributed Recoverability}

\begin{definition}
\begin{equation}
R^{dist}(\mathbf{X}) = \frac{1}{|H|} \sum_{i \in H} R_i(\tilde{\mathbf{x}}_i)
\end{equation}
\end{definition}

\begin{definition}
\begin{equation}
R_i^{min} = \min_{i \in H} R_i(\tilde{\mathbf{x}}_i)
\end{equation}
\end{definition}

\begin{proposition}
\begin{equation}
\min_i R_i \le R^{dist} \le \max_i R_i
\end{equation}
\end{proposition}

%%%%%%%%%%%%%%%%%%%%%%%%%%%%%%%%%%%%%%%%%%%%%%%%%%%%%%%%%%%%%%%%%%%%%%%%%%%%%%%%
\section{Failure Modes}

\begin{itemize}
\item \textbf{Local violation:} Node attempts inadmissible action.
\item \textbf{Consensus lag:} Delayed propagation of state.
\item \textbf{Partition:} Graph disconnects.
\end{itemize}

The kernel mitigates these through validation, consensus, and recovery.

%%%%%%%%%%%%%%%%%%%%%%%%%%%%%%%%%%%%%%%%%%%%%%%%%%%%%%%%%%%%%%%%%%%%%%%%%%%%%%%%
\section{Interception Example}

In a five-node network, a node attempts to increase autonomy beyond capacity. Validation blocks execution, consensus propagates the violation, and recovery dynamics restore system balance.

%%%%%%%%%%%%%%%%%%%%%%%%%%%%%%%%%%%%%%%%%%%%%%%%%%%%%%%%%%%%%%%%%%%%%%%%%%%%%%%%
\section{Discussion}

The Trust Kernel shows that admissibility and recoverability can be enforced in distributed systems. The architecture composes local invariants into global guarantees under standard fault assumptions.

%%%%%%%%%%%%%%%%%%%%%%%%%%%%%%%%%%%%%%%%%%%%%%%%%%%%%%%%%%%%%%%%%%%%%%%%%%%%%%%%
\section{Conclusion}

We presented a distributed architecture for governed autonomy. The Trust Kernel integrates validation, attestation, and recovery to enforce capacity constraints across nodes.

\end{document}
